% Part: first-order-logic
% Chapter: syntax-and-semantics
% Section: unique-readability

\documentclass[../../../include/open-logic-section]{subfiles}

\begin{document}

\olfileid{fol}{syn}{unq}

\olsection{Eenduidige leesbaarheid}

\begin{explain}
De wijze waarop we !!{formula}s hebben gedefinieerd, garandeert dat
iedere !!{formula} een \emph{eenduidige lezing} heeft. Er is dus in
wezen slechts één manier om haar volgens onze vormingsregels voor
!!{formula}s te construeren en slechts één manier om haar te
``interpreteren''. Anders zouden we ambigue !!{formula}s hebben:
!!{formula}s met meer dan één lezing of interpretatie. Dat willen we
uiteraard vermijden. Belangrijker is dat de meeste definities en
bewijzen die volgen zonder deze eigenschap niet zouden werken.

Dit wordt wellicht het duidelijkst wanneer we nagaan wat er zou
gebeuren bij gebrekkige vormingsregels voor !!{formula}s die geen
eenduidige leesbaarheid garanderen. We hadden bijvoorbeeld de haakjes
in de vormingsregels voor connectieven kunnen vergeten en het volgende
kunnen toelaten:
\begin{quote}
Als $!A$ en $!B$ !!{formula}s zijn, dan is $!A \lif !B$ dat ook.
\end{quote}
Uitgaande van een atomaire formule $!D$ zouden we dan $!D \lif !D$
kunnen vormen. Daaruit zouden we samen met $!D$ de formule $!D \lif !D
\lif !D$ verkrijgen. Dat kan echter op twee manieren:
\begin{enumerate}
\item We nemen $!D$ als $!A$ en $!D \lif !D$ als $!B$.
\item We nemen $!A$ als $!D \lif !D$ en $!B$ als~$!D$.  
\end{enumerate}
Dienovereenkomstig zijn er twee manieren om de !!{formula}~$!D \lif !D
\lif !D$ te ``lezen''. Zij heeft de vorm $!B \lif !C$, waarbij $!B$
gelijk is aan $!D$ en $!C$ aan $!D \lif !D$, maar zij heeft
\emph{eveneens} de vorm $!B \lif !C$, waarbij $!B$ gelijk is aan $!D
\lif !D$ en $!C$ aan~$!D$.

In zo'n geval werken onze definities niet altijd. Bij de definitie van
de !!{main operator} van een formule zeggen we bijvoorbeeld dat in een
formule van de vorm $!B \lif !C$ het aangegeven voorkomen van~$\lif$
de !!{main operator} is. Als we de formule $!D \lif !D \lif !D$ echter
op de twee bovengenoemde manieren met $!B \lif !C$ kunnen vergelijken,
dan is in het ene geval het eerste voorkomen van $\lif$ de !!{main
operator} en in het andere geval het tweede. De !!{main operator} moet
echter een \emph{functie} van de !!{formula} zijn: iedere !!{formula}
moet precies één voorkomen van een !!{main operator} hebben.
\end{explain}

\begin{lem}
Het aantal linker- en rechterhaakjes in !!a{formula}~$!A$ is gelijk.
\end{lem}

\begin{proof}
We bewijzen dit met inductie over de constructie van $!A$. Daarvoor
zijn twee zaken nodig: (a) eerst bewijzen we dat alle atomaire formules
de bedoelde eigenschap hebben (de inductiebasis); (b) vervolgens
bewijzen we dat uit reeds gegeven formules met deze eigenschap alleen
nieuwe formules met dezelfde eigenschap worden geconstrueerd.

Laat $l(!A)$ het aantal linkerhaakjes en $r(!A)$ het aantal
rechterhaakjes in~$!A$ zijn. Evenzo geven $l(t)$ en $r(t)$ het aantal
linker- respectievelijk rechterhaakjes in een term~$t$ aan.

\begin{prob}
    Bewijs dat voor iedere term~$t$ geldt: $l(t) = r(t)$.
\end{prob}

\begin{enumerate}
\tagitem{prvFalse}{\indcase{!A}{\lfalse}{$\indfrm$ heeft $0$ linker- en $0$
    rechterhaakjes.}}{}

\tagitem{prvTrue}{\indcase{!A}{\ltrue}{$\indfrm$ heeft $0$ linker- en $0$
    rechterhaakjes.}}{}

\item \indcase{!A}{\Atom{R}{t_1,\dots,t_n}}{$l(\indfrm) = 1 + l(t_1) +
  \dots + l(t_n) = 1 + r(t_1) + \dots + r(t_n) = r(\indfrm)$. Hier
  gebruiken we het als oefening overgelaten feit dat $l(t) = r(t)$
  voor iedere term~$t$.}

\item \indcase{!A}{\eq[t_1][t_2]}{$l(\indfrm) = l(t_1) + l(t_2) =
  r(t_1) + r(t_2) = r(\indfrm)$.}

\tagitem{prvNot}{\indcase{!A}{\lnot !B}{Volgens de inductiehypothese
    geldt $l(!B) = r(!B)$. Dus $l(\indfrm) = l(!B) = r(!B) =
    r(\indfrm)$.}}{}

\item \indcase{!A}{(!B \ast !C)}{Volgens de inductiehypothese geldt $l(!B) =
  r(!B)$ en $l(!C) = r(!C)$. Dus $l(\indfrm) = 1 + l(!B) + l(!C) =
  1 + r(!B) + r(!C) = r(\indfrm)$.}

\tagitem{prvAll}{\indcase{!A}{\lforall[x][!B]}{Volgens de
    inductiehypothese geldt $l(!B) = r(!B)$. Dus $l(\indfrm) = l(!B) = r(!B) =
    r(\indfrm)$.}}{}

% Print case for \lexists if prvEx and not prvAll
\tagitem{prvAll,notprvEx}{}{\indcase{!A}{\lexists[x][!B]}{Volgens de
    inductiehypothese geldt $l(!B) = r(!B)$. Dus $l(\indfrm) = l(!B) = r(!B) =
    r(\indfrm)$.}}

% Just say similarly if prvEx and prvAll
\tagitem{notprvAll,notprvEx}{}{\indcase{!A}{\lexists[x][!B]}{Analoog.}}
\end{enumerate}
\end{proof}

\begin{defn}[Echte prefix]
Een tekenreeks $!B$ is een \emph{echte prefix} van een tekenreeks~$!A$
als de concatenatie van $!B$ met een niet-lege tekenreeks gelijk is
aan~$!A$.
\end{defn}

\begin{lem}\ollabel{lem:no-prefix}
Als $!A$ !!a{formula} is en $!B$ een echte prefix van $!A$, dan is
$!B$ geen !!a{formula}.
\end{lem}

\begin{proof}
Oefening.
\end{proof}

\begin{prob}
Bewijs \olref[fol][syn][unq]{lem:no-prefix}.
\end{prob}

\begin{prop}
\ollabel{prop:unique-atomic}
Als $!A$ een atomaire !!{formula} is, voldoet zij aan precies één van
de volgende voorwaarden.
\begin{enumerate}
\tagitem{prvFalse}{$!A \ident \lfalse$.}{}
\tagitem{prvTrue}{$!A \ident \ltrue$.}{}
\item $!A \ident \Atom{R}{t_1,\dots,t_n}$, waarbij $R$ een $n$-plaatsig
  !!{predicate} is, $t_1$, \dots, $t_n$ termen zijn en elk van $R$,
  $t_1$, \dots, $t_n$ eenduidig bepaald is.
\item $!A \ident \eq[t_1][t_2]$, waarbij $t_1$ en $t_2$ eenduidig
  bepaalde termen zijn.
\end{enumerate}
\end{prop}

\begin{proof}
Oefening.
\end{proof}

\begin{prob}
Bewijs \olref[fol][syn][unq]{prop:unique-atomic} (Aanwijzing: formuleer
en bewijs een versie van \olref[fol][syn][unq]{lem:no-prefix} voor termen.)
\end{prob}

\begin{prop}[Eenduidige leesbaarheid]
Iedere !!{formula} voldoet aan precies één van de volgende voorwaarden.
\begin{enumerate}
\item $!A$ is atomair.

\tagitem{prvNot}{$!A$ heeft de vorm $\lnot !B$.}{}

\tagitem{prvAnd}{$!A$ heeft de vorm $(!B \land !C)$.}{}

\tagitem{prvOr}{$!A$ heeft de vorm $(!B \lor !C)$.}{}

\tagitem{prvIf}{$!A$ heeft de vorm $(!B \lif !C)$.}{}

\tagitem{prvIff}{$!A$ heeft de vorm $(!B \liff !C)$.}{}

\tagitem{prvAll}{$!A$ heeft de vorm $\lforall[x][!B]$.}{}

\tagitem{prvEx}{$!A$ heeft de vorm $\lexists[x][!B]$.}{}
\end{enumerate}
Bovendien zijn in ieder geval $!B$, of $!B$ en $!C$, eenduidig
bepaald. Er bestaan dus bijvoorbeeld geen verschillende paren $!B$,
$!C$ en $!B'$, $!C'$ waarvoor $!A$ zowel de vorm
\iftag{prvIf}{$(!B \lif !C)$ als $(!B' \lif !C')$ heeft.}{
    \iftag{prvOr}{$(!B \lor !C)$ als $(!B' \lor !C')$ heeft.}{
      \iftag{prvAnd}{$(!B \land !C)$ als $(!B' \land !C')$ heeft.}{}}}
\end{prop}

\begin{proof}
Volgens de vormingsregels moet iedere niet-atomaire !!{formula}
beginnen met een openingshaakje~(, \iftag{prvNot}{$\lnot$,}{} of een
kwantor. Omgekeerd moet iedere !!{formula} die met een van de volgende
symbolen begint atomair zijn: !!a{predicate}, !!a{function},
!!a{constant}\iftag{prvFalse}{, $\lfalse$}{}\iftag{prvTrue}{, $\ltrue$}{}.

We hoeven dus alleen nog aan te tonen dat als $!A$ zowel de vorm $(!B
\ast !C)$ als de vorm $(!B' \mathbin{\ast'} !C')$ heeft, dan $!B
\ident !B'$, $!C \ident !C'$ en $\ast = {\ast'}$.

Veronderstel dus zowel $!A \ident (!B \ast !C)$ als $!A \ident (!B'
\mathbin{\ast'} !C')$. Dan geldt $!B \ident !B'$ of niet. In het eerste
geval volgen duidelijk $\ast = {\ast'}$ en $!C \ident !C'$, want het
zijn dan deelreeksen van $!A$ die op dezelfde plaats beginnen en
dezelfde lengte hebben. Het andere geval is $!B \not\ident !B'$. Omdat
$!B$ en $!B'$ beide deelreeksen van $!A$ zijn die op dezelfde plaats
beginnen, moet de ene een echte prefix van de andere zijn. Dit is
echter onmogelijk volgens \olref{lem:no-prefix}.
\end{proof}
\end{document}
